\documentclass[12pt]{article}
\usepackage[utf8]{inputenc}
\usepackage[margin=1in]{geometry}   % 设置页边距
\usepackage{amsmath, amssymb}       % 数学符号与公式
\usepackage{enumitem}               % 自定义列表
\usepackage{unicode-math}
\usepackage[english]{babel}         % 设置整个文档为英文


% 标题与作者信息（可自行修改）
\title{Solution for Chapter 3}
\author{Zimo Guo}
\date{}

\begin{document}

\maketitle

% 使用 \noindent 去掉段落缩进，让题号更清晰
\noindent \textbf{3.3}

\begin{enumerate}[label=(\alph*), left=0pt] % left=0pt 使子题与题号对齐
    \item 
    From left to right, if $\{\mathcal{X}^{+} \cup \{\textbf{x}_{m+1}\}, \mathcal{X}^{-}\}$ is 
    linearly separable, then the hyperplane is located between the negative samples and the 
    point $\textbf{x}_{m+1}$. If $\{\mathcal{X}^{+}, \{\textbf{x}_{m+1}\} \cup \mathcal{X}^{-}\}$ 
    is linearly separable, then the hyperplane is located in the positive samples and the point 
    $\textbf{x}_{m+1}$. So the hyperplane is the line going though origin and $\textbf{x}_{m+1}$.

    From right to left, if $\{\mathcal{X}^{+}, \mathcal{X}^{-}\}$ 
    is linearly separable by a hyperplane going through the origin and $\textbf{x}_{m+1}$, then 
    there exist some hyperplane located between $\textbf{x}_{m+1}$ and the negative samples and going 
    though the origin can separate $\{\mathcal{X}^{+} \cup \{\textbf{x}_{m+1}\}\}$ and 
    $\{\mathcal{X}^{-}\}$, and some hyperplane located between $\textbf{x}_{m+1}$ and the 
    positive samples and going through the origin can separate $\{\mathcal{X}^{+}\}$ and 
    $\{\{\textbf{x}_{m+1}\} \cup \mathcal{X}^{-}\}$. So they are linearly separable.

    \item 
    \begin{align*}
        \text{Right hand side} &= 2\sum_{k=0}^{d-1} \binom{m-1}{k} + 2\sum_{k=0}^{d-2} \binom{m-1}{k} \\
        &= 2\sum_{k=0}^{d-1} \frac{m-k}{m}\binom{m}{k} + 2\sum_{k=0}^{d-2} \frac{m-k}{m}\binom{m}{k} \\
        &= 2[\frac{m-0}{m}\binom{m}{0} + \frac{m-1}{m}\binom{m}{1} + \frac{m-2}{m}\binom{m}{2} + \cdots + 
        \frac{m-(d-1)}{m}\binom{m}{d-1} \\
        &+ \frac{m-0}{m}\binom{m}{0} + \frac{m-1}{m}\binom{m}{1} + \frac{m-2}{m}\binom{m}{2} + \cdots + 
        \frac{m-(d-2)}{m}\binom{m}{d-2}] \\
    \end{align*}
    
    Compare with $C(m+1, d)$, we need to prove the following equality,
    \begin{align*}
        \frac{d-1}{m}\binom{m}{d-1} &= \sum_{k=0}^{d-2}\frac{m-2k}{m}\binom{m}{k} \\
        \iff (d-1)\binom{m}{d-1} &= (m-2k)\sum_{k=0}^{d-2}\binom{m}{k}
    \end{align*}

    \item 
    $\mathcal{\Psi}(x)$ is the vector of $f(x)$ and is the new input in the linear function. There 
    exists m samples such that every p-subset of $\{\mathcal{\Psi}(x_{1}), \mathcal{\Psi}(x_{2}), 
    \dots, \mathcal{\Psi}(x_{m})\}$ is linearly independent, so $\textbf{a}$ and $\mathcal{\Psi}(x)$
    play same role with $\textbf{w}$ and $\textbf{x}$ in the formula in previous case. Use the equality
    proved in (b), and obviously $C(m, d) > C(m, d-1)$, so the maxmium number of dichotomy is $C(m, p)$. 
    That is $\prod_{\mathcal{F}} (m) = C(m, p) = 2 \sum_{i=0}^{p-1} \binom{m-1}{i}$.
\end{enumerate}

\vspace{1em}  % 题目之间的垂直间距

\noindent \textbf{3.4}

\vspace{1em}  % 题目之间的垂直间距

\noindent \textbf{3.5}

$\prod (\mathcal{G}, S)$ is the number of ways for $\mathcal{G}$ to label sample set $S$. So

\begin{align*}
    \Pi (\mathcal{G}, S) &= |\{(g(x_{1}), g(x_{2}), \dots, g(x_{m})), g \in \mathcal{G}\}| \\
\end{align*}.

We want to bound Rademancher complexity in terms of $\prod (\mathcal{G}, S)$, so we need to represent
it in terms of the elements in $(g(x_{1}), g(x_{2}), \dots, g(x_{m}))$. Let $(g_{j}(x_{1}), g_{j}(x_{2}), 
\dots, g_{j}(x_{m}))$ be $d_{j}$, and use Massart's lemma, 

\begin{align*}
    \mathfrak{R}_{m}(\mathcal{G}) &= E_{S}[E_{\sigma}[\frac{\sup\limits_{g \in \mathcal{G}} \sum_{i=1}^{m} \sigma_{i} g(x_{i})}{m}]] \\
    &= E_{S}[E_{\sigma}[\frac{1}{m} \sup\limits_{d \in \textbf{b}} \boldsymbol{\sigma} \mathbf{d}]] \\
    &\le E_{S}[\sqrt{\frac{2 \log \Pi (\mathcal{G}, S)}{m}}]
\end{align*}

\vspace{1em}  % 题目之间的垂直间距

\noindent \textbf{3.6}

\begin{enumerate}[label=(\alph*), left=0pt] 
    \item 
    \begin{align*}
        \mathfrak{R}_{m}(\mathcal{G}) &= E_{S}[E_{\sigma}[\frac{\sup \sum \limits_{i=1}^{m} 
            \sigma_{i} g(x_{i})}{m}]] \\
        &= E_{S}[E_{\sigma}[\sigma \mathbf{g}]] \\
        &= E_{S}[E_{\sigma}[\sigma] E_{\sigma} [\mathbf{g}]]
    \end{align*}

    Because $\sigma$ is sequence of random variables valued -1 or 1, its expectation is zero
    vector. So $\mathfrak{R}_{m}(\mathcal{G}) = 0$ for singleton hypothesis set.

    \item

\end{enumerate}

\vspace{1em}  % 题目之间的垂直间距

\noindent \textbf{3.8}

\begin{enumerate}[label=(\alph*), left=0pt]
    \item 
    \begin{align*}
        \mathfrak{R}_{m}(\mathcal{H}) &= E_{S}[E_{\sigma}[\frac{\sup \limits_{\alpha h \in \alpha\mathcal{H}}
            \sum \limits_{i=1}^{m} \sigma_{i} h(x_{i})}{m}]] \\
        &= E_{S}[E_{\sigma}[\frac{|\alpha| \sup \limits_{\alpha h \in \alpha\mathcal{H}}
            \sum \limits_{i=1}^{m} \sigma_{i} h(x_{i})}{m}]] \\
        &= E_{S}[|\alpha| E_{\sigma}[\frac{\sup \limits_{\alpha h \in \alpha\mathcal{H}}
            \sum \limits_{i=1}^{m} \sigma_{i} h(x_{i})}{m}]] \\
        &= |\alpha| E_{S}[E_{\sigma}[\frac{\sup \limits_{\alpha h \in \alpha\mathcal{H}}
            \sum \limits_{i=1}^{m} \sigma_{i} h(x_{i})}{m}]] \\
        &= |\alpha| \mathfrak{R}_{m}(\mathcal{H})
    \end{align*}

    \item 
    \begin{align*}
        \mathfrak{R}_{m}(\mathcal{H} + \mathcal{H}^{'}) &= E_{S}[E_{\sigma}[\frac{\sup \limits_{h \in \mathcal{H},
        h^{'} \in \mathcal{H}^{'}} \sum \limits_{i=1}^{m} \sigma_{i} [h(x_{i}) + h^{'}(x_{i})]}{m}]] \\
        &= E_{S}[E_{\sigma}[\frac{\sup \limits_{h \in \mathcal{H}}\sum \limits_{i=1}^{m} \sigma_{i} h(x_{i})}{m}]] 
        + E_{S}[E_{\sigma}[\frac{\sup \limits_{h^{'} \in \mathcal{H^{'}}} \sum \limits_{i=1}^{m} \sigma_{i} h^{'}(x_{i})}{m}]] \\
        &= \mathfrak{R}_{m}(\mathcal{H}) + \mathfrak{R}_{m}(\mathcal{H}^{'})
    \end{align*}

    \item 
    If $|f(x_{1}) - f(x_{2})| \le L |x_{1} - x_{2}|$, $f(x)$ is called \textbf{L-Lipschitz function}.

    To use Talagrand's lemma, we need to prove that $max(a, b)$ is L-Lipschitz function. 
    \begin{align*}
        |max(a_{1}, b_{1}) - max(a_{2}, b_{2})| &= \frac{1}{2} |(a_{1} - a_{2}) + (b_{1} - b_{2}) + 
        (|a_{1} - b_{1}| - |a_{2} - b_{2}|)| \\
        &\le \frac{1}{2} ||a_{1} - a_{2}| + |b_{1} - b_{2}| + ||a_{1} - b_{1}| - |a_{2} - b_{2}||| \\
        &\le \frac{1}{2} ||a_{1} - a_{2}| + |b_{1} - b_{2}| + |(a_{1} - b_{1}) - (a_{2} - b_{2})|| \\
        &= \frac{1}{2} ||a_{1} - a_{2}| + |b_{1} - b_{2}| + |(a_{1} - a_{2}) - (b_{1} - b_{2})|| \\
        &\le \frac{1}{2} ||a_{1} - a_{2}| + |b_{1} - b_{2}| + |a_{1} - a_{2}| - |b_{1} - b_{2}|| \\
        &= ||a_{1} - a_{2}| + |b_{1} - b_{2}|| \\
        &= ||(a_{1}, b_{1}) - (a_{2}, b_{2})|| _{1}
    \end{align*}

    So $max(a, b)$ is 1-Lipschitz function. Then we can use Talagrand's lemma, 
    \begin{align*}
        \mathfrak{R}_{m}({\max(h, h^{'})})
    \end{align*}
\end{enumerate}

\vspace{1em}  % 题目之间的垂直间距

\noindent \textbf{3.9}

$ h = \max(h_{1} + h_{2} - 1, 0)$, So 
\begin{align*}
    \hat{\mathfrak{R}}_{S}(\mathcal{H})
    &= E[\sup \frac{\sum \sigma_{i} h(x_{i})}{m}] \\
    &= E[\sup \frac{\sum \sigma_{i} \max(h_{1}(x_{i}) + h_{2}(x_{i}) - 1, 0)}{m}] \\
    &\le E[\sup \frac{\sum \sigma_{i} (h_{1}(x_{i}) + h_{2}(x_{i}) - 1)}{m}] \\ 
    &= E[\sup \frac{\sum \sigma_{i} h_{1}(x_{i})}{m}] + E[\sup \frac{\sum \sigma_{i} h_{2}(x_{i})}{m}] + 
        E[\sup \frac{\sum \sigma_{i}}{m}] \\
    &= \hat{\mathfrak{R}}_{S}(\mathcal{H}_{1}) + \hat{\mathfrak{R}}_{S}(\mathcal{H}_{2})
\end{align*}

\vspace{1em}  % 题目之间的垂直间距

\noindent \textbf{3.11}

\begin{enumerate}[label=(\alph*), left=0pt]
    \item 
    \begin{align*}
        \hat{\mathfrak{R}}_{S}(\mathcal{H}) &= E[\sup \frac{\sum \limits_{i=1}^{m} \sigma_{i} \sum \limits_{j=1}^{n_{2}}
            w_{j}\sigma (u_{j} x_{i})}{m}] \\
        &= E[\sup \frac{\sum \limits_{j=1}^{n_{2}} w_{j} \sum \limits_{i=1}^{m} \sigma_{i}
            \sigma(u_{j} x_{i})}{m}] \\
        &= \frac{\Lambda^{'}}{m} E[\sup |\sum \limits_{i=1}^{m} \sigma_{i} \sigma(\mathbf{u} x_{i})|]
    \end{align*}

    \item 
    \begin{align*}
        \hat{\mathfrak{R}}_{S}(\mathcal{H}) &= \frac{\Lambda^{'}}{m} E[\sup |\sum \limits_{i=1}^{m} \sigma_{i} \sigma(\mathbf{u}
            x_{i})|] \\
        &\le \frac{\Lambda^{'}}{m} E[\sup |\sum \limits_{i=1}^{m} \sigma_{i} \mathbf{u}x_{i}|] \\
        &\le \frac{\Lambda^{'}}{m} E[\sup \sum \limits_{i=1}^{m} |\sigma_{i} \mathbf{u}x_{i}|] \\
        &= \frac{\Lambda^{'}}{m} E[\sup \sum \limits_{i=1}^{m} |\mathbf{u}| |\sigma_{i} x_{i}|] \\
        &\le \frac{\Lambda^{'} \Lambda}{m} E[\sum \limits_{i=1}^{m} \sigma_{i} x_{i}] \\
        &= \Lambda^{'} \hat{\mathfrak{R}}_{S}(\mathcal{H^{'}})
    \end{align*}

    \item 
    Cauchy-Schwarz inequality is following,
    \begin{align*}
        ||u v||_{2} \le ||u||_{2} ||v||_{2}
    \end{align*}.

    Use it in $\hat{\mathfrak{R}}_{S}(\mathcal{H^{'}})$, 
    \begin{align*}
        \hat{\mathfrak{R}}_{S}(\mathcal{H^{'}}) &= \frac{1}{m} E[\sup \sum \limits_{i=1}^{m} \sigma_{1} s\mathbf{u}x_{i}] \\
        &= \frac{1}{m} E[\sup \sum \limits_{i=1}^{m} \sigma_{1} \mathbf{u}x_{i}] \\
        &= \frac{\Lambda}{m} E[||\sum \limits_{i=1}^{m} \sigma_{i} x_{i}||_{2}]
    \end{align*}

    \item
    \begin{align*}
        \hat{\mathfrak{R}}_{S}(\mathcal{H^{'}}) &= \frac{\Lambda}{m} E[||\sum \limits_{i=1}^{m} \sigma_{i}x_{i}||_{2}] \\
        &\le \frac{\Lambda}{m} \sqrt{E[||\sum \limits_{i=1}^{m} \sigma_{i}x_{i}||_{2}^{2}]}
    \end{align*}

    \item
    Use the previous results and the inequality (3.14) in the book,
    \begin{align*}
        \mathfrak{R}_{m}(\mathcal{H}) &\le \hat{\mathfrak{R}}_{S}(\mathcal{H}) + \sqrt{\frac{\log \frac{2}{\delta}}{2m}} \\
        &\le \frac{\Lambda^{'} \Lambda}{m} \sqrt{E[m r^{2}]} + \sqrt{\frac{\log \frac{2}{\delta}}{2m}}
    \end{align*}
\end{enumerate}

\vspace{1em}  % 题目之间的垂直间距

\noindent \textbf{3.12}

\vspace{1em}  % 题目之间的垂直间距

\noindent \textbf{3.14}

For a hypothesis set $\mathcal{H}$, the possible most labeling results is $|\mathfrak{H}|$, 
where every $h$ in the set give a different predict value. So in binary classification tasks,
$2^{m} \le |\mathcal{H}|$. Get $m \le \log_{2} |\mathcal{H}|$.

In another jiaodu, if VC-dimention of $\mathcal{H}$ is $\log_{2} |\mathcal{H}| + 1$, then the 
hypothesis set should can give $2^{\log_{2} |\mathcal{H}|} = 2|\mathcal{H}|$ dichotomy, which 
is impossible.

\vspace{1em}  % 题目之间的垂直间距

\noindent \textbf{3.19}

If the dimention of $\mathcal{F}$ is $r$, consider a set has $m = r+1$ points, one hypothesis
in $\mathcal{H}$ is linear dependent, that is for $x_{r+1}$, whether it is in this hypothesis
could be judged from other $x$. So if VC-dimention of $\mathcal{H}$ is $r+1$ and $x_{r+1}$ is
linear dependent with $x_{p}$ and $x_{q}$, $2^{r+1}$ dichotomy is impossible. So $d < r$.



\end{document}